\documentclass[12pt, a4paper]{article}

% Packages
\usepackage[margin=1in]{geometry}
\usepackage{booktabs}
\usepackage{hyperref}
\usepackage{enumitem}
\usepackage{setspace}
\usepackage[numbers]{natbib}

% Formatting
\onehalfspacing

\hypersetup{
    colorlinks=true,
    linkcolor=blue,
    citecolor=blue,
    urlcolor=blue
}

\begin{document}

% ============ HEADER ============
\begin{center}
    {\large \textbf{National Institute of Technology Calicut}}\\[0.2cm]
    {\normalsize Department of Computer Science \& Engineering}\\[0.8cm]
    
    {\Large \textbf{Proposed Research Topic}}\\[0.4cm]
    
    {\large \textbf{Multimodal Anomaly Detection in Industrial Quality Inspection\\ using Vision-Language Models}}\\[0.8cm]
    
    \begin{tabular}{rl}
        \textbf{Student:} & [Your Name] \\
        \textbf{Roll No:} & [Your Roll Number] \\
        \textbf{Guide:} & Prof.\ [Professor Name] \\
        \textbf{Date:} & February 13, 2026 \\
    \end{tabular}
\end{center}

\vspace{0.5cm}
\hrule
\vspace{0.5cm}

% ============ 1. AREA OF INTEREST ============
\section{Area of Interest}

\textbf{Multimodal anomaly detection} involves using multiple types of data (e.g., images and text) together to identify items that deviate from expected normal patterns. In industrial manufacturing, this translates to automated quality inspection --- detecting product defects such as scratches, dents, missing components, or structural damage on production lines.

Current anomaly detection methods predominantly rely on \textbf{single-modality} (image-only) approaches and require training on normal samples for every new product category, which limits their scalability. Recent advances in \textbf{Vision-Language Models} (VLMs) offer a promising alternative by jointly understanding images and natural language, potentially enabling \textbf{zero-shot} anomaly detection --- detecting defects without any product-specific training.

% ============ 2. INITIAL FINDINGS ============
\section{Initial Findings from Literature}

I studied several recent works in this area. Here are the key findings from three representative papers:

\subsection{Paper 1: WinCLIP \cite{jeong2023winclip}}
\begin{itemize}[leftmargin=*]
    \item \textbf{What they did:} Applied the CLIP vision-language model for zero-shot anomaly detection in manufacturing. Used text prompts like ``a photo of a damaged [product]'' vs ``a photo of a normal [product]'' and a sliding window approach for defect localization.
    \item \textbf{Key result:} Achieved 91.8\% image-level AUROC on the MVTec AD benchmark \textbf{without any training}, competitive with many supervised methods.
    \item \textbf{Limitation:} Relies on hand-crafted prompts; cannot explain \textit{why} something is anomalous.
\end{itemize}

\subsection{Paper 2: AnomalyGPT \cite{gu2024anomalygpt}}
\begin{itemize}[leftmargin=*]
    \item \textbf{What they did:} Used a Large Vision-Language Model (LVLM) fine-tuned on simulated anomaly data for industrial defect detection. The model can detect anomalies, localize them, and describe them in natural language through dialogue.
    \item \textbf{Key result:} Achieves state-of-the-art results on MVTec AD with the added ability to \textit{explain} detected defects in words.
    \item \textbf{Limitation:} Requires fine-tuning on simulated data; not truly zero-shot.
\end{itemize}

\subsection{Paper 3: MMAD Benchmark \cite{jiang2025mmad}}
\begin{itemize}[leftmargin=*]
    \item \textbf{What they did:} Created a comprehensive benchmark to evaluate Multimodal LLMs (GPT-4o, Claude, open-source models) on industrial anomaly detection across 7 subtasks.
    \item \textbf{Key result:} Even the best models (GPT-4o) show \textbf{significant room for improvement}, indicating this is an open and active research problem.
    \item \textbf{Takeaway:} There is a clear opportunity to develop better MLLM-based anomaly detection approaches.
\end{itemize}

% ============ 3. IDENTIFIED GAP ============
\section{Preliminary Research Gap}

From these papers, I observe:
\begin{enumerate}[leftmargin=*]
    \item CLIP-based methods (WinCLIP) work well for zero-shot scoring but cannot explain anomalies.
    \item MLLM-based methods (AnomalyGPT) can explain anomalies but require fine-tuning.
    \item Current MLLMs (as shown by MMAD benchmark) still perform poorly out-of-the-box for industrial anomaly detection.
\end{enumerate}

\textbf{Gap:} A systematic study of how to effectively use Vision-Language Models and open-source MLLMs (e.g., LLaVA) together for zero-shot anomaly detection --- combining quantitative scoring with natural language reasoning --- is currently lacking.

% ============ 4. PROPOSED DIRECTION ============
\section{Proposed Research Direction}

I would like to investigate:

\begin{quote}
\textit{``How can pre-trained Vision-Language Models and Multimodal Large Language Models be effectively leveraged for zero-shot anomaly detection in industrial quality inspection, achieving robust defect detection and interpretable explanations without task-specific training?''}
\end{quote}

\textbf{Broadly, this would involve:}
\begin{itemize}[leftmargin=*]
    \item Using CLIP for anomaly scoring with structured prompt engineering
    \item Using an open-source MLLM (LLaVA) for defect reasoning and explanation
    \item Evaluating on standard benchmarks (MVTec AD dataset, 15 product categories)
    \item Comparing against traditional supervised baselines
\end{itemize}

% ============ 5. NEXT STEPS ============
\section{Next Steps}

\begin{itemize}[leftmargin=*]
    \item Conduct a more thorough literature survey (10--12 papers)
    \item Formalize the problem statement and methodology
    \item Set up the experimental environment and run preliminary experiments
    \item Prepare a detailed proposal for mid-semester review
\end{itemize}

% ============ REFERENCES ============
\begin{thebibliography}{3}

\bibitem{jeong2023winclip}
J.~Jeong, Y.~Zou, T.~Kim, D.~Zhang, A.~Ravichandran, and O.~Dabeer.
\newblock WinCLIP: Zero-/few-shot anomaly classification and segmentation.
\newblock In \emph{CVPR}, 2023.

\bibitem{gu2024anomalygpt}
Z.~Gu et~al.
\newblock AnomalyGPT: Detecting industrial anomalies using large vision-language models.
\newblock In \emph{AAAI}, 2024.

\bibitem{jiang2025mmad}
X.~Jiang et~al.
\newblock MMAD: Multi-modal anomaly detection benchmark.
\newblock In \emph{ICLR}, 2025.

\end{thebibliography}

\end{document}
